\section{Activation}
\label{sec:activation}

\subsection{Spec freeze and activation timeline}

\begin{longtable}{p{3.5cm}p{10cm}}
\toprule
\textbf{Date} & \textbf{Milestone} \\
\midrule
\endhead
2025-12-15 & \Quasar{} 3.0 spec freeze (\textsc{lp-020}). Per-LLM chain spec frozen alongside. \\
2025-12-16 -- 2025-12-22 & Genesis dry-runs on testnet for the \texttt{zen4} family. Validators rehearse \BlockSTM{} ordering, A-Chain bridge handshake, and \Lux{} Cloud federation. \\
2025-12-23 -- 2025-12-24 & Final genesis state generated. Validator set seeded. Operator pool stake-locked. \\
2025-12-25 & \Quasar{} 3.0 activation. \texttt{zen4} per-LLM chains launch with checkpoints derived from upstream HF weights. First on-chain training epoch begins. \\
\bottomrule
\end{longtable}

\subsection{Genesis state}

For each \Zen{} model in \S\ref{sec:zen}, genesis state contains:
\begin{itemize}[nosep]
\item The HF source commit pinned as the initial weights root.
\item The architecture artifact for the upstream foundation, registered as research artifact zero.
\item An empty receipts tree (no training has occurred on chain yet).
\item An empty attribution tree (initial weights are imported, not earned on chain).
\item The reward token contract deployed and the DAO treasury allocated.
\end{itemize}

A genesis import is itself a transaction signed by the \Zoo{} multisig that holds upstream license attestations. Subsequent training accrues attribution to on-chain contributors as defined in \S\ref{sec:training}.

\subsection{Operator onboarding}

Operators wishing to commit capacity at activation must:
\begin{enumerate}[nosep]
\item Register an account on \Lux{} Cloud (\textsc{lp-136}).
\item Submit hardware claim (\GPU{} class, count, geo, network).
\item Stake-lock per the chain's $\kappa$ multiplier.
\item Pass an attestation handshake against A-Chain (\textsc{lp-134}).
\end{enumerate}

The onboarding window closes 24\,h before activation; post-activation onboarding is continuous but new operators wait one epoch before receiving shard claims (this is a slashing-safety buffer, not a centralization choice).

\subsection{Post-activation roadmap}

The post-activation roadmap is sequenced as four phases beyond 2025-12-25; exact dates are governed by on-chain milestones rather than calendar quarters.

\begin{itemize}[nosep]
\item \textbf{Phase I (post-activation)}: \texttt{zen4} identity LoRA fine-tunes are first-on-chain training runs; results validate the protocol against existing centralized fine-tunes.
\item \textbf{Phase II}: \texttt{zen5} chains genesis. \texttt{zen5} introduces MoDE (Mixture of Diverse Experts) and is the first \Zen{} generation trained natively on chain rather than imported from upstream.
\item \textbf{Phase III}: cross-chain inference revenue routing from \Hanzo{} AI Chain to \Zoo{} per-LLM chains live in production.
\item \textbf{Phase IV}: F-Chain (\textsc{lp-013}) FHE inference live for \texttt{zen-omni} and \texttt{zen-vl-*} chains.
\end{itemize}

\subsection{Conclusion}

Per-LLM chains are the smallest design that binds training data, compute, research, and revenue into a single audit-eligible ledger. They are not a generic blockchain hosting model files; they are an execution environment in which the lifecycle of a single model is the chain's reason to exist. The \Zen{} family demonstrates the design at scale: eighteen models, eighteen sovereign chains, one settlement layer. \Quasar{} 3.0 (\textsc{lp-020}) provides the security; the \Quasar{}-Native App Stack (\textsc{lp-133}) provides the substrate; \Quasar{}\GPU{} (\textsc{lp-132}) provides scheduling; A-Chain (\textsc{lp-134}) provides attestation; \Lux{} Cloud (\textsc{lp-136}) provides capacity federation. The \Zoo{} per-LLM chains compose these into the open foundation-model lab that centralized labs cannot match.

Activation: 2025-12-25.
